

\chapter{\label{III_3} Généraliser le processus} 

L'éditeur logiciel aurait un rôle à jouer dans la mise à disposition de nouvelles fonctionnalités innovantes auprès d'institutions patrimoniales, d'autant plus lorsqu'il s'agit de collections audiovisuelles. En effet, selon Jonas Arnold, 

\begin{quote}
	Les institutions de mémoire contenant des contenus audiovisuels sont prédestinées à explorer les avantages du \gls{ml} ou de l'intelligence hybride pour la mise en valeur et l'accès à leurs contenus audiovisuels. [...] Dans ce cas, la collaboration avec des partenaires, tels que des prestataires de services ou des portails, s'impose.” \footcite{zotero-item-584}. 
\end{quote}

Or, le processus décrit en seconde partie, centré sur une validation utilisateur, a été analysé à partir d'un corpus spécifique, choisit comme exemple de mise en application. Si le rôle de l'éditeur est celui de la mise à disposition de fonctionnalités innovantes pour le domaine patrimonial, cellec-ci doivent pouvoir être généralisées à tous les tyoes de corpus audiovisuels hébergés par le logiciel. Ce chapitre propose ainsi une mise en perspective des propositions élaborées au cours de notre étude.




\section{\label{III_3_A}Statut des données : données automatiques et données manuelles/métiers}

L'introduction de processus fondés sur l'intelligence artificielle au sein d'un logiciel de gestion des collections modifie la nature et le statut des données produites sur les documents. Comme le souligne Jonas Arnold dans le \textit{Livre blanc sur l'apprentissage automatique dans les archives}, les données automatiquement produites possèdent une nature "fluide" puisqu'elles sont amenées à évoluer au rythme des améliorations des outils, internes ou externes \footcite{zotero-item-584}. 

(transition avec sous partie suivante : nécessité de généraliser le processus pour que les données ne bougent pas trop) (transition ausis avec refonte de S-Movies : le contrôle qualité exige une meilleure navigation de l'interface te une vue d'ensemble)

La nouvelle nature des données implique la mise en place d'une intelligence hybride qui impliquerait par exemple que l'utilisateur métier supervise la cohérence entre les entités du modèle. Ainsi, tandis que le processus automatique nourrit des entités inférieures, celles du fragment et de l'item, il reviendrait à l'utilisateur de garantir les liens avec les entités supérieures, celle de l'Oeuvre, et la cohérence du reversement des données au sein du modèle dans son ensemble. Le reversement des données automatique en \textit{backend} prévu au sein de la seconde partie est un revesrement théorique, qu'il revient à l'utilisateur de vérifier la cohérence. 
Afin d'assurer la vérification de la cohérence du modèle de données, deux dispositifs pourront être mis en place au sein du logiciel. 
D'une part, afin de pouvoir garantir l'homogénéité de l'indexation de ces données, le logiciel doit prévoir leur traçabilité via une "historisation ou un contrôle des versions" \footcite{zotero-item-584}. 
Skinsoft prévoit des champs de métadonnées de gestion au sein du logiciel contenant par exemple la date de modification de la notice, le dernier utilisateur à l'avoir modifié. Les métadonnées peuvent être définies comme de l'information structurée décrivant des données \footcite{colavizza2026}. Le même type de métadonnées pourra être conservé au sein de champs dédiés aux métadonnées de processus, telles que les outils utilisés, les paramètres appliqués, l'utilisateur ayant validé l'indexation, la date. Ces métadonnées apparaitraient dans une notice de Fragment par exemple. 
La traçablité des métadonnées produites permettrait ainsi à l'archiviste non seulement de suivre les évolutions et de veiller à la conformité de leur structure, mais également aux données d'être réemployées à des fins de réentraînement des machines. Ainsi, si le processus est modifié par l'emploi d'un nouvel outil par exemple, les métadonnées de processus permettront une réadaptation efficace sans modifier l'interface utilisateur encore davantage. 
D'autre part, pour une question de transparence et de visibilité au sein de l'interface, l'origine des données doit être excplicitement signalée. Si des métadonnées de processus sont intégrées au sein des notices, les données issues de processus automatiques doivent être clairement identifiée d'emblée au sein des grilles de résultats de recherche, par exemple. Cette identification peut prendre la forme d'un code couleur ou de tags de provenance apposés sur les données. Ce système permettrait à l'utilisateur de vérifier immédiatement si les données sont associées au bonne entités du modèle FRBR. 
Ainsi, la mise en place d'une intelligence hybride consiste à permettre la production automatique de données patrimoniales, à la condition que l'utilisateur garantisse leur cohérence et leur validation au sein du modèle de données. 
afin de gérer les erreurs potentielles et les bugs : système de tickets et de tests pour régler ensuite les régressions de l'application. 
\section{\label{III_3_B}Qualité des données : une tâche à part entière}

Ainsi, le contrôle de la cohérence des données au sein du modèle de données s'impose comme une tâche à part entière, à laquelle s'ajoute le contrôle de la qualité des données elles-mêmes. Si un pipeline de validation permet d'encadrer la création des données descriptives, les erreurs peuvent persister et dépendent du temps dédié à la validation. Dès lors, l'évaluation de la qualité des données produites s'ajoute à la vérification de leur cohérence structurelle au sein du modèle FRBR. 

À ce stade du processus, la qualité des données ne désigne pas de métrique technique mais plutôt une valeur métier, en lien avec les normes descriptives en vigueur. Deux niveaux d'exigence de qualité peuvent être établis. Premièrement, les données produites doivent être évaluées en lien avec les normes de description employées au sein de l'espace de travail. Par exemple, S-Movies est fondé sur les règles de description des archives (\gls{rda}) conformément aux reccommandations de la \gls{fiaf}. Deuxièmement, les données pourront être vérifiées en accord avec les règles internes et propres à chaque institution patrimoniale. Par exemple, les thésaurus, les autorités, les dates. Comme nous l'avions mentionné, les mots-clés produits sur le corpus étudié en seconde partie sont contrôlés par un thésaurus "keyword". Il s'agira alors de vérifier si les mots-clés versés dans l'application existent bien sous une forme contrôlée au sein d'un thésaurus dédié. 

Toutefois, ces vérifications laissées à l'utilisateur métier peuvent prendre un temps considérable et ne permettent pas de couvrir l'entiereté de la masse de données produites automatiquement. Dans ce cadre, l'usage de grands modèles de langage (\gls{llm}) peut être envisagé afin de contrôler la qualité des données au niveau institutionnel. En effet, Skinsoft développe déjà l'intégration d'un \gls{llm} pour les fonctionnalités de recherche en langage naturel et de \gls{rag}, entraîné sur les règles d'indexation métier générales et spécifique aux institutions. Le même \gls{llm} pourra être utilisé afin de pointer les incohérences au sein des notices contenant des données indexées automatiquement et manuellement. Un tel outil permettrait également d'homogénéiser l'ensemble des données indexées. Ainsi, l'utilisation d'un outil basé sur l'IA supplémentaire permettrait d'améliorer la production de données homogènes en regard des normes en vigueur, contrairement aux données indexées manuellement, souvent hétérogènes.  

Utiliser des LLM pour corriger les données et exécuter ces tâches de validation d'une manière hybride entre l'humaine et la machine, l'IA ne dégrade pas forcément la qualité des métadonnées mais peut l'améliorer en produisant des métadonnées homogènes au regard des normes. Selon Nikki Wise, certaines études montrent même que les utilisateurs préfèrent "les métadonnées générées par des machines ou par des systèmes hybrides à celles créées par des humains, en invoquant leur contexte plus riche, leur spécificité et leur exhaustivité [Traduction personnelle]" \footnote{\textit{"human users preferred machine-and hybrid-generated metadata over human-generated content, citing its increasedcontext, specificity, and comprehensiveness"} \cite{wise2026}}." Les capacités d'interprétation et d'analyse des \gls{llm} s'avèrent particulièrement efficaces pour assister l'utilisateur, et ce, même avec un entraînement limité, offrant ainsi "une perspective prometteuse pour les flux de travail hybrides associant humains et IA dans la création et la correction des métadonnées. [Traduction personnelle]" \footnote{\textit{"suggest promise for hybrid human-AI workflows in metadata creation and remediation"} \cite{wise2026}}.
Ce contrôle de la qualité des données ne peut toutefois pas s'effectuer sans leur traçabilité, prévue par les métadonnées de processus décrites précédemment, permettant d'avoir accès à l'historique des modifications et au contexte de création de ces données.

%Contrôle qualité et apprentissage continu : erreurs identifiées comme données d'annotation 
%Enfin, cette traçabilité alimente une boucle d'apprentissage continu. Les erreurs détectées par l'archiviste lors du contrôle qualité ne sont pas seulement corrigées dans la notice : elles sont requalifiées en données d'annotation. Réinjectées dans le pipeline, ces corrections enrichissent le jeu d'entraînement des modèles, perfectionnant ainsi les performances des futures propositions automatiques. 
\section{\label{III_3_C} Une automatisation nécesaire ?}

\subsection{La réalité des besoins}
Au delà des contraintes posées par la segmentation et l'indexation automatisées, se pose la question de la pertinence du déploiement de telles fonctionnalités, en lien avec les besoins professionnels actuels. En effet, un éditeur tel que Skinsoft est animé par une logique commerciale d'innovation qui contraste avec les habitudes métier de gestion des collections. Au sein d'un \textit{guide pour une IA responsable dans le secteur culturel}, le Ministère de la Culture rappelle que dans le cadre d'une mise en place de fonctionnalités fondées sur l'IA, il s'agit "d’examiner le besoin à couvrir et de vérifier la capacité des systèmes d’IA à satisfaire ce besoin dans des conditions optimales" ; le cas échéant, de favoriser des alternatives moins consommatrices que l’IA si elles sont suffisantes à répondre au besoin." \footcite{2025a}.
Dans le cadre de notre étude, notre proposition se limite à l'expérimentation d'un processus à partir d'un besoin précis identifié sur un corpus spécifique. Or, ce besoin fait partie d'une réalité globale documentaire, celle de la nécessité d'exploiter davantage le contenu des archives audiovisuelles, dont le caractère temporel rend les caractéristiques inexploitées, inadaptées aux outils de recherche existants. Comme l'explique le chercheur Bruno Bachimont, 
\begin{quote} 
	La temporalité des séquences audiovisuelles a pour conséquence qu’il est impossible de consulter rapidement un document, de le "feuilleter" pour trouver un élément que l’on recherche : pour feuilleter une heure de vidéo, il faut une heure. \footcite{bachimont}. 
\end{quote} 

L'automatisation permettrait alors d'agir comme accélérateur d'une découvrabilité fine des contenus audiovisuels et d'une accélération de leur numérisation, par le temps réduit de leur indexation. Dans ce contexte,  l'éditeur de logiciel occupe un rôle d'intermédiaire technique entre le développement innovant de l'IA et les besoins patrimoniaux. À terme, les institutions culturelles occupent le rôle d'intermédiaire du contenu ainsi nouvellement indexé auprès du public. 

\subsection{Le coût matériel et humain de la supervision humaine}
Toutefois, nous avons pu observer au cours de notre stage certaines institutions refuser le projet d'implémentation de l'IA au sein du logiciel, pour des raisons de perte de transparence sur le traitement des données et de réticences quant à la confidentialité des données exploitées. Notre proposition est précisément axée sur un processus régit par une transparence maximale, permise par la validation indispensable du processus par l'utilisateur. En effet, le processus proposé dans notre étude, fondé sur la \gls{Vision par ordinateur} via des méthodes d'\gls{appsup}, est une alternative pertinente pour instaurer la transparence. Les outils que nous avons proposé permettent une liberté de configuration afin de conserver le contrôle sur les tâches automatisées. Les risques associés habituellement aux modèles d'IA tels que les risques d'hallucination, d'erreurs, de biais, seraient amoindris par le statut qui leur est donné : celui de facilitateur et d'assitant ponctuel, mais jamais celui de décisionnaire. 

Or, si ce processus centré sur la supervision humaine parait idéal, il impose des contraintes matérielles et financières importantes. D'une part, l'éditeur devra superviser le développement du processus via l'annotation de larges jeux de données, pour entraîner un modèle comme \gls{yolo}. D'autre part, l'institution cliente devra mettre en place des ressources dédiées à la validation du processus et au contrôle qualité des données automatiquement produites. L'implémentation de la segmentation et de l'indexation automatiques, devra ainsi s'effectuer au cas par cas, selon les besoins identifiés par les retours clients. Nous pouvons dès lors nous demander si l'implémentation du processus idéal tel que nous l'avons décrit, constitue un gain de temps de travail ou une charge supplémentaire.



Afin de satisfaire un besoin plus global, il est justement pertinent de mettre le modèle de données au premier plan de l'implémentation, afin de rapeller la capacité d'un modèle d'IA à s'adapter à des pratiques métier.
 

%Vers une publication conforme sur le web 
%éditeur comme vecteur de du web de données liées : préparer les données à leur publication sur lweb via des référentiels et une structuration frbr 

\subsection{Proposer un module paramétrable}
Afin de faire face à la disparité des besoins et des moyens entre les institutions, l'implémentation d'une segmentation et d'une indexation automatiques pourra s'effectuer sous la forme d'un module paramétrable et optionnel. En effet, étant donné la variabilité du modèle de données d'une institution à une autre, et des différences de configurations entre les espaces de travail, l'automatisation décrite pourra être intégrée à l'image du fonctionnement du logiciel, sous une forme paramétrable au cas par cas. De la même manière que l'éditeur propose la configurabilité de l'espace de travail par des masques, méta masques et menus, un dialogue en amont pourra être établi afin de définir les modalités d'implantation de la fonctionnalité. La modification du modèle de données par l'ajout d'une nouvelle entité pouvant recueillir les séquences pourra notamment être proposée comme une option. De plus, l'éditeur pourra prévoir un environnement logiciel de test afin de garantir une appropriation des fonctionnalités et d'en modifier le paramétrage au cours des pratiques. 
Ainsi, la segmentation et l'indexation automatiques pourront être fondées sur une architecture paramétrable, de la même manière que l'ensemble du logiciel. Le mode configurable de ces fonctionnalités par l'éditeur auront pour but d'une part de soulager les institutions de la mise en place de moyens couteux, et d'autre part de garantir la mise en place d'une IA comme outil d'assistance pour l'utilisateur.  




La nécessité d'implémenter une segmentation et une indexation automatiques est remise en perspective par la disparité des besoins patrimoniaux. Si la constitution de jeux de données d'entraînement patrimoniaux permettrait d'améliorer la qualité des prédictions, et la généralisation de l'automatisation à d'autres institutions clientes, la supervision humaine décrite au sein de notre processus implique des besoins et des coûts considérables. 
Ainsi, la pertinence d'une telle automatisation résiderait dans une implémentation configurable, adaptée aux besoins et aux moyens de chaque institution.


